\documentclass[12pt]{scrartcl}
\usepackage[hagavi]{azzam}

\title{Pre Screening OSK SMA (Intermediate)}
\author{150 menit}
\date{Jumat, 20 November 2026}

\begin{document}
\maketitle

\renewcommand*\contentsname{Daftar Isi}
\tableofcontents

\newpage

\noindent Tes ini terdiri dari 20 soal isian singkat:
5 soal aljabar, 5 soal teori bilangan, 5 soal kombinatorika, dan 5 soal geometri.
Tuliskan jawaban akhirnya saja. Setiap jawaban yang benar bernilai 1 poin,
sedangkan jawaban yang salah atau kosong bernilai 0. Dilarang menggunakan kalkulator.

\noindent Soal 1--2, 6--7, 11--12, dan 16--17 sama dengan paket Pre Screening
sebelumnya sebagai pemanasan. Selebihnya adalah soal baru yang diambil dari
AMC 12 dan AIME, dengan tingkat kesulitan sedikit di atas OSK.

\section{Aljabar}
\begin{enumerate}
	% Sama dengan paket 2026 10 16 no 1 (fungsi kuadrat)
	\item Diberikan fungsi kuadrat $f(x)=ax^{2}+bx+c$ yang memenuhi $f(3)=9$ dan $f(8)=64$.
	Jika $a \ne 1$, nilai dari $\dfrac{c-b}{a-1}$ adalah \dots

	% Sama dengan paket 2026 10 16 no 2 (barisan geometri)
	\item Misalkan $u_{1},u_{2},u_{3},\dots$ suatu barisan geometri dengan $u_{1}>u_{2}$.
	Jika $u_{2}=9$ dan $u_{5}+u_{7}=\dfrac{10}{3}u_{6}$, nilai dari $u_{1}$ adalah \dots

	% Diadaptasi dari AMC 12A 2023 no 10
	\item Bilangan real positif $x$ dan $y$ memenuhi $y^{3}=x^{2}$ dan $(y-x)^{2}=4y^{2}$.
	Nilai dari $x+y$ adalah \dots

	% Diadaptasi dari AIME I 2024 no 2
	\item Terdapat bilangan real $x$ dan $y$, keduanya lebih besar dari 1, sehingga
	$\log_{x}(y^{x})=\log_{y}(x^{4y})=10$. Nilai dari $xy$ adalah \dots

	% Diadaptasi dari AMC 12A 2024 no 15 (akar polinomial)
	\item Akar-akar dari $x^{3}+2x^{2}-x+3$ adalah $p$, $q$, dan $r$.
	Nilai dari $(p^{2}+4)(q^{2}+4)(r^{2}+4)$ adalah \dots
\end{enumerate}

\section{Teori Bilangan}
\begin{enumerate}[resume]
	% Sama dengan paket 2026 10 16 no 5 (bilangan kuadrat)
	\item Suatu bilangan empat digit $\overline{6ab9}$ merupakan suatu bilangan kuadrat.
	Nilai dari $a+b$ adalah \dots

	% Sama dengan paket 2026 10 16 no 6 (faktorial)
	\item Bilangan bulat positif terkecil $n$ sehingga $n!$ habis dibagi $1540$ adalah \dots

	% Diadaptasi dari AMC 12A 2023 no 4 (digit)
	\item Banyak digit dalam representasi desimal dari $8^{5}\cdot 5^{10}\cdot 15^{5}$
	adalah \dots

	% Diadaptasi dari AMC 12A 2024 no 9 (bilangan kuadrat)
	\item Misalkan $M$ bilangan bulat terbesar sehingga $M+1213$ dan $M+3773$
	keduanya merupakan bilangan kuadrat sempurna.
	Digit satuan dari $M$ adalah \dots

	% Diadaptasi dari AMC 12A 2024 no 17 (sistem persamaan bulat)
	\item Bilangan bulat $a$, $b$, dan $c$ memenuhi $ab+c=100$, $bc+a=87$,
	dan $ca+b=60$. Nilai dari $ab+bc+ca$ adalah \dots
\end{enumerate}

\section{Kombinatorika}
\begin{enumerate}[resume]
	% Sama dengan paket 2026 10 16 no 9 (kombinasi, jumlah digit)
	\item Bilangan KOTA adalah bilangan 4 angka yang tidak dimulai dengan angka 0 dan hasil
	penjumlahan semua digitnya sama dengan 9. Sebagai contoh, 2025 merupakan bilangan KOTA.
	Banyaknya bilangan KOTA adalah \dots

	% Sama dengan paket 2026 10 16 no 10 (prinsip sarang merpati)
	\item Afif menuliskan sepuluh bilangan bulat positif yang lebih kecil dari 20.
	Ia memastikan bahwa penjumlahan dua bilangan mana pun di antara sepuluh bilangan tersebut
	tidak sama dengan 20. Bilangan positif yang pasti ditulis Afif adalah \dots

	% Diadaptasi dari AIME I 2024 no 4 (peluang bersyarat)
	\item Jen mengikuti undian dengan memilih 4 bilangan berbeda dari
	$S=\{1,2,\dots,10\}$. Kemudian 4 bilangan dari $S$ diundi secara acak.
	Jen memenangkan hadiah jika setidaknya dua dari bilangan pilihannya keluar,
	dan memenangkan hadiah utama jika keempat bilangan pilihannya keluar.
	Peluang Jen memenangkan hadiah utama dengan syarat ia memenangkan hadiah
	dapat ditulis sebagai $\frac{m}{n}$ dengan $m$ dan $n$ relatif prima.
	Nilai dari $m+n$ adalah \dots

	% Diadaptasi dari AIME I 2024 no 6 (lintasan grid)
	\item Perhatikan lintasan sepanjang 16 langkah yang mengikuti garis grid dari sudut
	kiri bawah ke sudut kanan atas pada grid $8\times 8$.
	Banyak lintasan yang tepat empat kali berubah arah
	(dari horizontal ke vertikal atau sebaliknya) adalah \dots

	% Diadaptasi dari AIME I 2024 no 3 (permainan mengambil token)
	\item Alice dan Bob bermain permainan berikut. Terdapat setumpuk $n$ token.
	Keduanya bergiliran dengan Alice jalan pertama.
	Pada setiap giliran, pemain mengambil 1 token atau 4 token dari tumpukan.
	Pemain yang mengambil token terakhir menang.
	Banyak bilangan bulat positif $n\le 2024$ sehingga Bob memiliki strategi
	yang menjamin kemenangan, apa pun langkah Alice, adalah \dots
\end{enumerate}

\section{Geometri}
\begin{enumerate}[resume]
	% Sama dengan paket 2026 10 16 no 13 (segiempat)
	\item Sebuah persegi dibagi menjadi dua persegi panjang. Diketahui hasil penjumlahan kedua
	keliling persegi panjang tersebut adalah 66, maka luas persegi semula adalah \dots

	% Sama dengan paket 2026 10 16 no 14 (garis istimewa segitiga)
	\item Diberikan $ABC$ siku-siku sama kaki dengan panjang $BC=AB$ dan $L$ titik tengah $BC$.
	Titik $P$ pada sisi $AC$ sehingga $BP$ tegak lurus dengan $AL$.
	Jika panjang $CP=24\sqrt{2}$, panjang $AB$ adalah \dots

	% Diadaptasi dari AIME I 2024 no 5 (segiempat siklik)
	\item Persegi panjang $ABCD$ berukuran $AB=107$ dan $BC=16$,
	serta persegi panjang $EFGH$ berukuran $EF=184$ dan $FG=17$.
	Titik $D$, $E$, $C$, dan $F$ terletak pada satu garis dalam urutan tersebut,
	sedangkan $A$ dan $H$ terletak di sisi yang berlawanan terhadap garis itu.
	Jika titik $A$, $D$, $H$, dan $G$ terletak pada satu lingkaran,
	panjang $CE$ adalah \dots

	% Diadaptasi dari AMC 12A 2024 no 19 (segiempat siklik)
	\item Segiempat siklik $ABCD$ memiliki panjang $BC=CD=3$ dan $DA=5$
	dengan $\angle CDA=120^{\circ}$.
	Jika panjang diagonal yang lebih pendek dapat ditulis sebagai $\frac{m}{n}$
	dengan $m$ dan $n$ relatif prima, nilai dari $m+n$ adalah \dots

	% Diadaptasi dari AIME I 2024 no 10 (garis singgung lingkaran luar)
	\item Diberikan $\triangle ABC$ dengan panjang sisi $AB=5$, $BC=9$, $CA=10$.
	Garis singgung lingkaran luar $\triangle ABC$ di $B$ dan $C$ berpotongan di titik $D$,
	dan $\overline{AD}$ memotong lingkaran luar di $P\ne A$.
	Jika panjang $AP$ dapat ditulis sebagai $\frac{m}{n}$ dengan $m$ dan $n$ relatif prima,
	nilai dari $m+n$ adalah \dots
\end{enumerate}

% ============================================================
% SOLUSI SINGKAT (mulai dari halaman baru).
% Untuk menyembunyikan solusi, beri tanda % di awal setiap baris
% dari blok ini sampai tepat sebelum \end{document}.
% ============================================================
% % \newpage
% % \section*{Solusi Singkat}
% \begin{enumerate}
% 	\item Kurangkan kedua persamaan: $55a+5b=55$, sehingga $b=11-11a$.
% 	Dari $f(3)=9$ diperoleh $c=9-9a-3b=24a-24$.
% 	Jadi $c-b=35a-35$ dan $\dfrac{c-b}{a-1}=\boxed{35}$.
%
% 	\item Tulis $u_{5}+u_{7}=u_{6}\left(\dfrac{1}{r}+r\right)$ dengan $r$ rasio barisan,
% 	sehingga $r+\dfrac{1}{r}=\dfrac{10}{3}$. Diperoleh $r=3$ atau $r=\dfrac{1}{3}$.
% 	Karena $u_{1}>u_{2}$, maka $r=\dfrac{1}{3}$ dan $u_{1}=\dfrac{9}{1/3}=\boxed{27}$.
%
% 	\item Dari $(y-x)^{2}=4y^{2}$ diperoleh $|y-x|=2y$.
% 	Kasus $y-x=2y$ memberi $x=-y$, mustahil untuk $x,y$ positif.
% 	Jadi $y-x=-2y$, yaitu $x=3y$. Substitusi ke $y^{3}=x^{2}=9y^{2}$ memberi $y=9$ dan $x=27$.
% 	Nilai $x+y=\boxed{36}$.
%
% 	\item Dari $\log_{x}(y^{x})=x\log_{x}y=10$ diperoleh $\log_{x}y=\dfrac{10}{x}$.
% 	Dari $\log_{y}(x^{4y})=4y\log_{y}x=10$ diperoleh $\log_{y}x=\dfrac{5}{2y}$.
% 	Karena $(\log_{x}y)(\log_{y}x)=1$, maka $\dfrac{10}{x}\cdot\dfrac{5}{2y}=1$,
% 	sehingga $xy=\boxed{25}$.
%
% 	\item Tulis $P(x)=x^{3}+2x^{2}-x+3=(x-p)(x-q)(x-r)$.
% 	Maka $(p^{2}+4)(q^{2}+4)(r^{2}+4)=P(2i)P(-2i)$.
% 	Karena $P(2i)=-8i-8-2i+3=-5-10i$,
% 	hasilnya $|-5-10i|^{2}=25+100=\boxed{125}$.
%
% 	\item Akar kuadratnya terletak di antara 78 dan 83. Agar kuadratnya berakhiran 9,
% 	akarnya harus berakhiran 3 atau 7, jadi akarnya 83 dan bilangannya $83^{2}=6889$.
% 	Nilai $a+b=8+8=\boxed{16}$.
%
% 	\item Karena $1540=2^{2}\cdot 5\cdot 7\cdot 11$, $n!$ harus memuat faktor prima 11,
% 	jadi $n \ge 11$. Nilai $11!$ memuat $2^{2}$, $5$, $7$, dan $11$, jadi $n=\boxed{11}$.
%
% 	\item Tulis $8^{5}\cdot 5^{10}\cdot 15^{5}=2^{15}\cdot 5^{15}\cdot 3^{5}=243\cdot 10^{15}$.
% 	Bilangan ini terdiri dari tiga digit $243$ diikuti 15 nol,
% 	jadi banyak digitnya $\boxed{18}$.
%
% 	\item Tulis $M+1213=a^{2}$ dan $M+3773=b^{2}$, sehingga
% 	$(b-a)(b+a)=b^{2}-a^{2}=2560$.
% 	Agar $M=a^{2}-1213$ terbesar, $a$ harus terbesar, jadi $d=b-a$ diambil pembagi
% 	genap terkecil dari 2560, yaitu $d=2$ (nilai $d=1$ mustahil karena paritas).
% 	Diperoleh $b+a=1280$, $a=639$, dan $M=639^{2}-1213=407108$.
% 	Digit satuannya $\boxed{8}$.
%
% 	\item Kurangkan persamaan pertama dan kedua:
% 	$(a-c)(b-1)=13$. Jadi $b-1$ membagi 13 dan $b\in\{2,0,14,-12\}$.
% 	Hanya $b=-12$ yang memberi solusi bulat yang konsisten, yaitu $a=-9$ dan $c=-8$.
% 	Jadi $ab+bc+ca=108+96+72=\boxed{276}$.
%
% 	\item Tulis $a+b+c+d=9$ dengan $a \ge 1$. Misalkan $a'=a-1$, maka
% 	$a'+b+c+d=8$ dengan semua peubah $\ge 0$. Banyak solusinya
% 	$\binom{8+4-1}{4-1}=\binom{11}{3}=\boxed{165}$.
%
% 	\item Pasangan bilangan yang berjumlah 20 adalah $(1,19)$, $(2,18)$, \dots, $(9,11)$,
% 	serta $10+10$. Sepuluh bilangan yang dipilih tidak boleh memuat dua bilangan yang
% 	berpasangan, jadi tepat satu bilangan diambil dari tiap pasangan dan bilangan 10
% 	wajib ikut. Bilangan yang pasti ditulis adalah $\boxed{10}$.
%
% 	\item Banyak cara 4 bilangan keluar adalah $\binom{10}{4}=210$.
% 	Cara menang hadiah: $\binom{4}{2}\binom{6}{2}+\binom{4}{3}\binom{6}{1}+\binom{4}{4}
% 	=90+24+1=115$. Peluang bersyaratnya $1/115$, jadi $m+n=\boxed{116}$.
%
% 	\item Setiap lintasan terdiri dari 8 langkah kanan dan 8 langkah atas.
% 	Tepat 4 kali berubah arah berarti 5 blok selang-seling.
% 	Jika diawali langkah kanan, polanya K-A-K-A-K dengan 3 blok kanan dan 2 blok atas:
% 	$\binom{7}{2}\binom{7}{1}=21\cdot 7=147$ cara. Sama banyak jika diawali langkah atas.
% 	Totalnya $2\cdot 147=\boxed{294}$.
%
% 	\item Posisi kalah untuk pemain yang jalan adalah $n\equiv 0$ atau $2\pmod 5$.
% 	Sebab dari $5k$ atau $5k+2$ hanya bisa jalan ke posisi menang,
% 	sedangkan dari posisi lain selalu bisa jalan ke $5k$ atau $5k+2$.
% 	Bob (pemain kedua) menang tepat di posisi kalah untuk Alice.
% 	Banyak $n\le 2024$ yang kongruen 0 atau 2 modulo 5 adalah $\boxed{809}$.
%
% 	\item Misalkan sisi persegi $s$ dan potongannya sejajar salah satu sisi.
% 	Jumlah kedua keliling adalah $2(s+h_{1})+2(s+h_{2})=4s+2s=6s$ karena $h_{1}+h_{2}=s$.
% 	Dari $6s=66$ diperoleh $s=11$, jadi luasnya $\boxed{121}$.
%
% 	\item Letakkan $B=(0,0)$, $C=(s,0)$, $A=(0,s)$ dengan $s=AB$.
% 	Maka $L=(s/2,0)$, garis $BP$ adalah $y=x/2$, dan garis $AC$ adalah $x+y=s$,
% 	sehingga $P=(2s/3,s/3)$. Jadi $CP=s\sqrt{2}/3=24\sqrt{2}$, maka $s=\boxed{72}$.
%
% 	\item Letakkan $D=(0,0)$, $C=(107,0)$, $A=(0,16)$, $E=(e,0)$,
% 	$F=(f,0)$, $H=(e,-17)$, $G=(f,-17)$ dengan $f-e=184$.
% 	Lingkaran melalui $D$ dan $A$ berbentuk $x^{2}+y^{2}+ux-16y=0$.
% 	Substitusi $H$ dan $G$ memberi $u=-(e^{2}+561)/e=-(f^{2}+561)/f$,
% 	sehingga $ef=561$. Dari $e(e+184)=561$ diperoleh $e=3$
% 	(akar negatifnya ditolak karena $E$ di antara $D$ dan $C$).
% 	Jadi $CE=107-e=\boxed{104}$.
%
% 	\item Aturan kosinus di $\triangle ADC$ memberi
% 	$AC^{2}=9+25-2\cdot 3\cdot 5\cos 120^{\circ}=49$, jadi $AC=7$.
% 	Karena segiempat siklik, $\angle ABC=60^{\circ}$, dan dari $\triangle ABC$
% 	diperoleh $AB=8$. Ptolemy memberi $7\cdot BD=8\cdot 3+3\cdot 5=39$,
% 	jadi $BD=39/7<7$. Nilai $m+n=39+7=\boxed{46}$.
%
% 	\item Garis $AD$ adalah simedian dari $A$, sehingga memotong $BC$ di $K$ dengan
% 	$BK:KC=AB^{2}:AC^{2}=1:4$. Diperoleh $BK=9/5$ dan $KC=36/5$.
% 	Stewart memberi $AK=26/5$. Dari kuasa titik $K$,
% 	$KP=KB\cdot KC/KA=162/65$.
% 	Jadi $AP=26/5+162/65=100/13$ dan $m+n=\boxed{113}$.
% \end{enumerate}

\end{document}
